%% ============================================================
%% Teil II: Robotik-Bibliotheken – NumPy, Matplotlib, OpenCV-Vorschau
%% Band 2 – Kapitel 6 bis 8
%% ============================================================

\part{Robotik-Bibliotheken}

% ─────────────────────────────────────────────────────────────
\chapter{Wir wollen Sensordaten filtern – NumPy für Robotik}
% ─────────────────────────────────────────────────────────────

\begin{lernziele}
Du kannst NumPy-Arrays für Sensor-Puffer einsetzen,
Mittelwerte und gleitende Fenster berechnen und weißt,
wann NumPy sich lohnt -- und wann einfaches Python reicht.
\end{lernziele}

\section{Situation}

Der HC-SR04 misst unruhig: 87, 91, 84, 88, 200 (Ausreißer!), 86, 89\,cm.
Ein einzelner Messwert ist unzuverlässig.
Für robuste Navigation braucht \texttt{navigation\_node}
gefilterte Werte -- Mittelwert, Median oder gleitendes Fenster.

\begin{aufgabeBox}
Implementiere einen Sensor-Puffer mit gleitendem Mittelwert
und Ausreißer-Filterung für den \texttt{arduino\_bridge\_node}.
Verwende NumPy für die Berechnung.
\end{aufgabeBox}

\section{Warum NumPy?}

\begin{wissenBox}
NumPy speichert Zahlen in kompakten C-Arrays -- nicht als Python-Objekte.
Operationen laufen als optimierter C-Code.
Für 1000 Messungen ist \texttt{np.mean(arr)} ca. 100-mal schneller
als eine Python-\texttt{for}-Schleife.

\begin{lstlisting}
import numpy as np

# Python-Schleife: langsam bei vielen Werten
messwerte = [87, 91, 84, 88, 86, 89]
mittelwert = sum(messwerte) / len(messwerte)   # 87.5

# NumPy: schnell, präzise, mehr Werkzeuge
arr = np.array(messwerte, dtype=np.float32)
np.mean(arr)    # 87.5
np.median(arr)  # 87.5 -- robust gegen Ausreißer
np.std(arr)     # Standardabweichung
np.clip(arr, 2, 400)  # Werte begrenzen: min=2cm, max=400cm
\end{lstlisting}

\textbf{Wann NumPy sich lohnt:} mehr als 100 Werte gleichzeitig,
Matrizenrechnung (Transformationen, Odometrie),
Bild-Arrays (OpenCV gibt numpy-Arrays zurück).

\textbf{Wann Python reicht:} einzelne Messwerte,
einfache Schwellenwertabfragen wie in \texttt{motor\_bridge.py}.
\end{wissenBox}

\subsection{Gleitender Mittelwert mit \texttt{deque} + NumPy}

\begin{lstlisting}
from collections import deque
import numpy as np

class SensorPuffer:
    """Gleitender Mittelwert über die letzten N Messungen."""

    def __init__(self, fenstergroesse: int = 5) -> None:
        self._puffer_rechts: deque[int] = deque(maxlen=fenstergroesse)
        self._puffer_links:  deque[int] = deque(maxlen=fenstergroesse)

    def hinzufuegen(self, rechts_cm: int, links_cm: int) -> None:
        self._puffer_rechts.append(rechts_cm)
        self._puffer_links.append(links_cm)

    def gefilterter_wert(self) -> tuple[float, float]:
        """Gibt Median beider Sensoren zurück (robust gegen Ausreißer)."""
        if not self._puffer_rechts:
            return 0.0, 0.0
        r = float(np.median(np.array(self._puffer_rechts, dtype=np.float32)))
        l = float(np.median(np.array(self._puffer_links,  dtype=np.float32)))
        return r, l

    def ist_bereit(self) -> bool:
        """Puffer muss mindestens halb voll sein."""
        return len(self._puffer_rechts) >= self._puffer_rechts.maxlen // 2
\end{lstlisting}

\subsection{Integration in den Node}

\begin{lstlisting}
class ArduinoBridgeNode(Node):
    def __init__(self) -> None:
        super().__init__('arduino_bridge')
        self._puffer = SensorPuffer(fenstergroesse=5)
        self._timer = self.create_timer(0.1, self._lese_arduino)

    def _lese_arduino(self) -> None:
        daten = self._lese_serielle_zeile()
        if daten is None:
            return
        self._puffer.hinzufuegen(daten.rechts_cm, daten.links_cm)
        if self._puffer.ist_bereit():
            r, l = self._puffer.gefilterter_wert()
            # In Teil III: hier publishen
            self.get_logger().debug(f"Gefiltert: R={r:.1f}cm L={l:.1f}cm")
\end{lstlisting}

\begin{projektBox}[robotik-uno-q -- Projektstand]
\textbf{Heute:} \texttt{SensorPuffer} mit gleitendem Median.
Ausreißer wie "`200\,cm"' werden herausgefiltert.\\
\textbf{Nächstes Kapitel:} Sensordaten live visualisieren.
\end{projektBox}

\begin{merksatzBox}
\textbf{Merksatz:} \texttt{np.median()} ist robuster als \texttt{np.mean()} --
ein Ausreißer beeinflusst den Median kaum,
den Mittelwert dagegen stark.
\texttt{collections.deque(maxlen=N)} ist perfekt für gleitende Fenster:
alte Werte fallen automatisch heraus.
\end{merksatzBox}

\begin{fehlerBox}
\begin{itemize}
\item \texttt{np.array([87, "N/A"])} → \texttt{dtype=object},
  keine Rechenoperationen möglich.
  Validiere immer vor dem Einfügen in den Puffer.
\item NumPy für einzelne skalare Werte:
  \texttt{np.mean([87])} ist langsamer als \texttt{87.0} -- Overhead lohnt sich erst ab ~50 Werten.
\item \texttt{deque} ohne \texttt{maxlen}: wächst unbegrenzt -- Memory-Leak im Dauerbetrieb.
\end{itemize}
\end{fehlerBox}

\begin{uebungBox}
\begin{enumerate}
\item Erweitere \texttt{SensorPuffer} um eine Methode \texttt{ausreisser\_erkannt()},
  die \texttt{True} zurückgibt, wenn der letzte Wert mehr als 2 Standardabweichungen
  vom Mittelwert abweicht.
\item Miss die Ausführungszeit von \texttt{gefilterter\_wert()} mit \texttt{timeit}
  für Fenstergrößen von 5, 50 und 500.
\item Welche Fenstergröße macht bei 10\,Hz Sensor-Rate und 1\,Hz Navigation Sinn?
  Begründe.
\end{enumerate}
\end{uebungBox}

\begin{haubeBox}
\textbf{NumPy und die GIL}

Der Python Global Interpreter Lock (GIL) verhindert echte
Parallelität in Python-Threads. NumPy-Operationen laufen
jedoch als C-Code -- und C-Code gibt die GIL frei.
\texttt{np.median()} auf einem großen Array blockiert andere Threads nicht.
Das macht NumPy besonders wertvoll in ROS2-Nodes,
die Timer- und Subscriber-Callbacks gleichzeitig bedienen.
\end{haubeBox}

\begin{robotikBox}
In ROS2 liefert \texttt{sensor\_msgs/LaserScan} 360~oder mehr Entfernungswerte
als \texttt{float32[]}-Array -- direkt als NumPy-Array nutzbar.
Das gleiche Filtermuster (\texttt{np.median()}, \texttt{np.clip()})
funktioniert für Lidar, Kamera-Tiefeninformation und IMU-Daten.
\end{robotikBox}


% ─────────────────────────────────────────────────────────────
\chapter{Wir wollen Sensordaten visualisieren – Matplotlib live}
% ─────────────────────────────────────────────────────────────

\begin{lernziele}
Du kannst Sensor-Zeitreihen live plotten,
weißt warum rviz2 für ROS2-Nodes die bessere Alternative ist
und kennst den Einsatzbereich von Matplotlib in der Robotik.
\end{lernziele}

\section{Situation}

Beim Debuggen willst du sehen, wie die HC-SR04-Rohdaten
im Vergleich zum gefilterten Wert aussehen --
nicht als Zahlenkolonne im Terminal,
sondern als Kurve.

\begin{aufgabeBox}
Erstelle ein Standalone-Diagnosewerkzeug, das Rohdaten und gefilterte
Abstandswerte des HC-SR04 live plottet.
Das Tool läuft neben dem ROS2-Node, nicht als Node selbst.
\end{aufgabeBox}

\section{Matplotlib für Offline-Analyse}

\begin{lstlisting}
#!/usr/bin/env python3
# stufe2_ros2/tools/sensor_plotter.py
"""Live-Plot der HC-SR04-Rohdaten aus einer Log-CSV."""
import matplotlib.pyplot as plt
import matplotlib.animation as animation
import csv
import time
from collections import deque

MAX_PUNKTE = 100
zeitreihe:  deque[float] = deque(maxlen=MAX_PUNKTE)
rechts_roh: deque[float] = deque(maxlen=MAX_PUNKTE)
links_roh:  deque[float] = deque(maxlen=MAX_PUNKTE)

fig, ax = plt.subplots(figsize=(10, 4))
ax.set_title("HC-SR04 Abstandsdaten – robotik-uno-q")
ax.set_xlabel("Zeit (s)")
ax.set_ylabel("Abstand (cm)")
ax.set_ylim(0, 420)
ax.axhline(y=25, color='red',    linestyle='--', label='Stopp (25 cm)')
ax.axhline(y=50, color='orange', linestyle='--', label='Warn (50 cm)')

linie_r, = ax.plot([], [], 'b-',  label='Rechts (roh)',      linewidth=1)
linie_l, = ax.plot([], [], 'g-',  label='Links (roh)',       linewidth=1)
ax.legend(loc='upper right')

def aktualisiere(frame: int) -> list:
    """Liest neue Zeilen aus der Log-CSV und aktualisiert den Plot."""
    try:
        with open('logs/sensor.csv', 'r') as f:
            zeilen = list(csv.reader(f))[-MAX_PUNKTE:]
        for ts, r, l in zeilen:
            zeitreihe.append(float(ts))
            rechts_roh.append(float(r))
            links_roh.append(float(l))
    except (FileNotFoundError, ValueError):
        pass

    linie_r.set_data(range(len(rechts_roh)), list(rechts_roh))
    linie_l.set_data(range(len(links_roh)),  list(links_roh))
    ax.set_xlim(0, MAX_PUNKTE)
    return [linie_r, linie_l]

ani = animation.FuncAnimation(fig, aktualisiere, interval=200, blit=True)
plt.tight_layout()
plt.show()
\end{lstlisting}

\section{rviz2 -- die bessere Alternative im ROS2-Kontext}

\begin{haubeBox}
\textbf{Matplotlib vs. rviz2: wann was?}

Matplotlib eignet sich für:
\begin{itemize}
\item Offline-Analyse von aufgezeichneten Daten
\item Einfache Zeitreihen außerhalb von ROS2
\item Berichte, Präsentationen, Paper
\end{itemize}

rviz2 eignet sich für:
\begin{itemize}
\item Live-Visualisierung von ROS2-Topics
\item 3D-Koordinatensysteme, Robotermodelle, Laserscans
\item Wegpunkte und Navigation
\end{itemize}

\textbf{Faustregel:} Wenn der Roboter läuft: rviz2.
Wenn du Daten nachträglichanalysierst: Matplotlib.
\end{haubeBox}

\begin{lstlisting}[language=bash]
# rviz2 starten (im ROS2-Container)
source /opt/ros/jazzy/setup.bash
rviz2

# In rviz2: Add → By topic → /sensor/rechts → Range
# Zeigt den Sensor live als 3D-Kegel
\end{lstlisting}

\begin{projektBox}[robotik-uno-q -- Projektstand]
\textbf{Heute:} \texttt{tools/sensor\_plotter.py} als Diagnosewerkzeug.\\
\textbf{Nächstes Kapitel:} Kamera-Anbindung -- Vorschau auf Band~5.
\end{projektBox}

\begin{merksatzBox}
\textbf{Merksatz:} Matplotlib für Offline-Analyse, rviz2 für Live-Daten in ROS2.
Beide Werkzeuge haben ihren Platz -- das \texttt{tools/}-Verzeichnis
ist genau für solche Hilfswerkzeuge gedacht.
\end{merksatzBox}


% ─────────────────────────────────────────────────────────────
\chapter{Vorschau: Kamera als Sensor – OpenCV und Band~5}
% ─────────────────────────────────────────────────────────────

\begin{lernziele}
Du weißt, wie eine Kamera in Python angesprochen wird,
verstehst den Zusammenhang zwischen OpenCV-Arrays
und ROS2-Bildnachrichten und weißt, was in Band~5 folgt.
\end{lernziele}

\section{Situation}

Ein Ultraschallsensor sieht nur geradeaus.
Für autonome Navigation -- Schilder lesen, Objekte erkennen,
Personen verfolgen -- braucht der Roboter eine Kamera.

\begin{aufgabeBox}
Lese ein einzelnes Bild von einer USB-Kamera,
zeige es an und verstehe, wie ROS2 Kamerabilder überträgt.
Diese Einheit ist eine Vorschau -- die vollständige Bildverarbeitung
folgt in Band~5 (Embedded AI).
\end{aufgabeBox}

\section{OpenCV: das Fundament}

\begin{lstlisting}
import cv2
import numpy as np

# Kamera öffnen (Index 0 = erste Kamera)
kamera = cv2.VideoCapture(0)

ret, bild = kamera.read()
# bild ist ein numpy-Array: shape = (480, 640, 3) -- H x W x BGR

# Graustufen
grau = cv2.cvtColor(bild, cv2.COLOR_BGR2GRAY)

# Anzeigen
cv2.imshow("Kamera", bild)
cv2.waitKey(0)

kamera.release()
cv2.destroyAllWindows()
\end{lstlisting}

\section{ROS2 und Kamerabilder}

\begin{wissenBox}
ROS2 überträgt Kamerabilder als \texttt{sensor\_msgs/Image}.
Das \texttt{cv\_bridge}-Paket konvertiert zwischen
OpenCV-numpy-Arrays und ROS2-Nachrichten:

\begin{lstlisting}
from cv_bridge import CvBridge
import cv2

bridge = CvBridge()

# numpy -> ROS2
ros_bild = bridge.cv2_to_imgmsg(bild, encoding="bgr8")

# ROS2 -> numpy
bild = bridge.imgmsg_to_cv2(ros_bild, desired_encoding="bgr8")
\end{lstlisting}

In Band~5 kommt YOLO dazu: Das Erkennungsmodell läuft auf dem
numpy-Array und gibt Bounding-Boxes zurück.
\end{wissenBox}

\begin{projektBox}[robotik-uno-q -- Projektstand]
\textbf{Heute:} Kamera-Grundkenntnisse als Vorschau.\\
\textbf{Teil~III:} rclpy -- der eigentliche ROS2-Einstieg beginnt jetzt.
\end{projektBox}

\begin{robotikBox}
Band~5 (Embedded AI) baut direkt auf diesem Kapitel auf:
USB-Kamera → OpenCV → YOLO-Erkennung → ROS2-Topic \texttt{/kamera/erkennung}.
Der Roboter wird dann nicht nur Hindernissen ausweichen,
sondern auch Personen, Schilder und Objekte erkennen.
\end{robotikBox}
